\section*{Item 2}\label{it:2}

\textbf{Enunciado}: Um conversor Forward possui os seguintes parâmetros: $V_i = 125$ V, $V_o = 50$ V e $R = 25~\Omega$. A frequência de comutação é de $250$ kHz. Determine:

\medskip
\textbf{a)} a relação de espiras do transformador $N_1/N_2$ tal que o ciclo de trabalho seja $0,3$;

\medskip
\textbf{Resolução}: para o conversor Forward operando em modo contínuo, a característica de conversão é dada por~\cite{forward_notas}
\begin{equation*}
	\frac{V_o}{V_i}
	=
	\frac{N_2}{N_1}D.
\end{equation*}

Isolando a relação de espiras,
\begin{equation*}
	\frac{N_1}{N_2}
	=
	\frac{D V_i}{V_o}.
\end{equation*}

Substituindo os valores fornecidos,
\begin{equation*}
	\frac{N_1}{N_2}
	=
	\frac{0{,}3\cdot125}{50}
	=
	0{,}75.
\end{equation*}

\medskip
\textbf{b)} a indutância $L$ (no circuito de saída), de modo que a corrente mínima em $L$ seja $40\%$ da corrente média;

\medskip
\textbf{Resolução}: inicialmente, calcula-se a corrente média na saída,
\begin{equation*}
	I_{L,\mathrm{med}}
	=
	\frac{V_o}{R}
	=
	\frac{50}{25}
	=
	2~\mathrm{A}.
\end{equation*}

Como a corrente mínima deve ser $40\%$ da corrente média,
\begin{equation*}
	I_{L,\min}
	=
	0{,}4I_{L,\mathrm{med}}
	=
	0{,}8~\mathrm{A}.
\end{equation*}

A ondulação de corrente no indutor é dada por~\cite{forward_notas}
\begin{equation*}
	\Delta i_L
	=
	\frac{V_o}{L}(1-D)T_s.
\end{equation*}

Como
\begin{equation*}
	I_{L,\min}
	=
	I_{L,\mathrm{med}}
	-
	\frac{\Delta i_L}{2},
\end{equation*}
obtém-se
\begin{equation*}
	0{,}4I_{L,\mathrm{med}}
	=
	I_{L,\mathrm{med}}
	-
	\frac{V_o(1-D)T_s}{2L}.
\end{equation*}

Isolando a indutância,
\begin{equation*}
	L
	.
\end{equation*}

Logo,
\begin{equation*}
	\begin{aligned}
		L
		&=
		\frac{V_o(1-D)4\cdot10^{-6}}
		{2(1-0{,}4)I_{L,\mathrm{med}}}\\
		&=\frac{50(1-0{,}3)4\cdot10^{-6}}
		{2(0{,}6)(2)}\\
		&=
		58{,}3~\mu\mathrm{H}.
	\end{aligned}
\end{equation*}

\medskip
\textbf{c)} a capacitância necessária para limitar a ondulação (ripple) da tensão de saída a $0{,}5\%$.

\medskip
\textbf{Resolução}: a ondulação relativa da tensão de saída para um conversor Forward é dada por~\cite{forward_notas}
\begin{equation*}
	\frac{\Delta V_o}{V_o}
	=
	\frac{1-D}{8LCf_s^2D}.
\end{equation*}

Deseja-se limitar a ondulação a
\begin{equation*}
	\frac{\Delta V_o}{V_o}=0{,}5\%=0{,}005.
\end{equation*}

Isolando a capacitância,
\begin{equation*}
	C
	=
	\frac{1-D}
	{8L f_s^2D
		\left(\frac{\Delta V_o}{V_o}\right)}.
\end{equation*}

Substituindo os valores obtidos,
\begin{equation*}
	\begin{aligned}
		C
		&=
		\frac{1-0{,}3}
		{8(58{,}3\cdot10^{-6})
			(250\cdot10^3)^2(0{,}3)(0{,}005)}
		\\
		&\approx4{,}8~\mu\mathrm{F}.
	\end{aligned}
\end{equation*}
